\chapter{Được và Mất}
\markboth{Được và Mất}{Gain and Loss}

\section{Một người được khen nhiều rồi sống mệt để giữ hình ảnh.}
\begin{truyen}{Một người được khen nhiều rồi sống mệt để giữ hình ảnh.}{Praise can turn into pressure.}
\chuhoa{C}{àng được khen giỏi, anh càng không dám hở ra chỗ yếu. Sai một tí là cuống lên, không hẳn vì lỗi to, mà vì sợ tụt khỏi cái hình tượng người ta đã trót gắn cho mình. Được tiếng thơm ban đầu nghe sướng, nhưng giữ nó lâu ngày nhiều khi mệt như vác đá.}
\textit{(The more he was praised as capable, the more he feared showing weakness. Every small mistake terrified him, not because the mistake was huge, but because others might admire him less. Some gains begin as joy, but if held too tightly they become burdens.)}
\end{truyen}
\begin{baihoc}
Được tiếng mà mất tự nhiên thì coi như lỗ nặng.
\textit{(Gaining a beautiful image while losing naturalness is no small price.)}
\end{baihoc}
\begin{ghinhoanh}
\textbf{Short English to remember:}

Praise can become pressure.
\end{ghinhoanh}
\begin{tuhocnhanh}
1. Em có đang sống để giữ một hình ảnh nào không? \textit{(Are you living to maintain an image?)}

2. Hình ảnh đó đang giúp em hay làm em mệt? \textit{(Is that image helping you or exhausting you?)}
\end{tuhocnhanh}
\ngancach

\section{Một học sinh mất một kỳ thi nhưng được một cách học đúng.}
\begin{truyen}{Một học sinh mất một kỳ thi nhưng được một cách học đúng.}{One failure can rescue your method.}
\chuhoa{T}{rượt một bài thi quan trọng, em sốc thật sự vì nghĩ bao công đổ sông đổ biển. Nhưng ngồi nhìn lại mới thấy bấy lâu mình học kiểu vá víu, thuộc để qua chứ không hiểu gốc. Rớt một lần thì đau, nhưng cũng nhờ vậy mới bỏ được kiểu học giả vờ chăm.}
\textit{(After failing an important exam, the student cried as if all effort had vanished. Yet that fall forced a review: the past method was mere memorization, not understanding. Failing one exam hurt, but it also opened a more sustainable way to learn.)}
\end{truyen}
\begin{baihoc}
Mất một điểm số mà nhặt lại được cách học đàng hoàng thì chưa chắc là thua.
\textit{(Losing one result may mean gaining a foundation.)}
\end{baihoc}
\begin{ghinhoanh}
\textbf{Short English to remember:}

A lost score can save your future.
\end{ghinhoanh}
\begin{tuhocnhanh}
1. Em có thất bại nào đã buộc em sửa cách làm chưa? \textit{(Has any failure forced you to change your method?)}

2. Em đang tiếc kết quả hay đang học từ nó? \textit{(Are you only grieving the result, or are you learning from it?)}
\end{tuhocnhanh}
\ngancach

\section{Một người được thêm tiền nhưng mất dần thời gian sống.}
\begin{truyen}{Một người được thêm tiền nhưng mất dần thời gian sống.}{More income can cost your inner life.}
\chuhoa{N}{hận thêm việc, anh khoái vì tiền vào đều hơn. Một thời gian sau thì bữa ăn nào cũng nuốt vội, sách bỏ đó, người nhà nói chuyện chưa tròn câu đã lại cắm đầu vào việc. Nhìn lại mới thấy ví có dày thêm chút, còn người thì xẹp xuống thấy rõ.}
\textit{(He took on more work and felt happy because income rose. Later he no longer read, rarely spoke with family, and always ate in a rush. Looking back, he realized he had become slightly richer but far emptier.)}
\end{truyen}
\begin{baihoc}
Không phải kiếm thêm được là sống thêm được.
\textit{(Not every gain makes life more full.)}
\end{baihoc}
\begin{ghinhoanh}
\textbf{Short English to remember:}

More money is not always more life.
\end{ghinhoanh}
\begin{tuhocnhanh}
1. Điều em đang theo đuổi đang lấy mất điều gì? \textit{(What is the thing you are chasing taking away from you?)}

2. Cái giá ấy có đáng không? \textit{(Is that price worth it?)}
\end{tuhocnhanh}
\ngancach

\section{Một người mất bạn vì nói thật đúng lúc.}
\begin{truyen}{Một người mất bạn vì nói thật đúng lúc.}{Truth can cost comfort.}
\chuhoa{K}{hi thấy bạn sắp làm bậy, cô không hùa theo cho êm. Cô nói thẳng nên bị giận, bị né, bị xem là nhiều chuyện. Nhưng lâu sau chính người bạn đó quay lại cảm ơn, vì giữa lúc ai cũng vuốt cho dễ chịu, ít ra có một người dám chặn mình lại.}
\textit{(When her friend was about to do something wrong, she did not stay silent to please. She spoke directly and was avoided for a while. Months later the friend returned with thanks, because in a crowd of easy approval, someone had dared to create discomfort in order to pull her back.)}
\end{truyen}
\begin{baihoc}
Có lúc mất lòng trước mắt còn đỡ hơn để người mình thương trượt dài.
\textit{(Some things lose short-term comfort in order to preserve long-term truth.)}
\end{baihoc}
\begin{ghinhoanh}
\textbf{Short English to remember:}

Truth may cost comfort.
\end{ghinhoanh}
\begin{tuhocnhanh}
1. Em thường chọn êm chuyện hay chọn điều đúng? \textit{(Do you usually choose comfort or what is right?)}

2. Có ai từng làm em khó chịu nhưng thực ra đang giúp em không? \textit{(Has anyone ever made you uncomfortable while actually helping you?)}
\end{tuhocnhanh}
\ngancach

\section{Một người chịu mất mặt để giữ lời hứa.}
\begin{truyen}{Một người chịu mất mặt để giữ lời hứa.}{Losing face can mean keeping value.}
\chuhoa{A}{nh từng hứa với học sinh rằng chấm sai thì xin lỗi công khai. Đến lúc chấm nhầm thật, anh biết mình hoàn toàn có thể lờ đi cho đỡ quê. Nhưng cuối giờ anh vẫn đứng ra nhận lỗi. Quê một lúc thôi, còn lòng tin thì giữ được lâu hơn.}
\textit{(He had promised his students that if he was wrong, he would apologize publicly. After grading incorrectly, he knew he could quietly let it pass, but at the end of class he admitted the mistake. He lost a little pride and gained something bigger: trust.)}
\end{truyen}
\begin{baihoc}
Có cái mất mặt bên ngoài lại cứu được cái mặt bên trong.
\textit{(Some outward losses grow inward value.)}
\end{baihoc}
\begin{ghinhoanh}
\textbf{Short English to remember:}

Humility builds trust.
\end{ghinhoanh}
\begin{tuhocnhanh}
1. Em có điều gì nên nhận sai thay vì tìm cách che? \textit{(Is there something you should admit instead of trying to hide?)}

2. Em muốn được nể vì vẻ ngoài hay được tin vì nhân cách? \textit{(Do you want to be admired for appearance or trusted for character?)}
\end{tuhocnhanh}
\ngancach
